%==========================================================================
\section{Introduction}\label{sec:introduction}

Statevector simulation stores all $2^n$ complex amplitudes of an $n$-qubit register and updates them gate by gate~\cite{nielsen2010quantum}. It remains the reference method for verifying circuits and algorithms of moderate width and arbitrary depth~\cite{preskill2018quantum,arute2019quantum}. Its cost is set by memory: thirty qubits in single precision (complex64) occupy 8~GB, and every gate reads and writes a large part of that state. On current hardware the arithmetic per amplitude is small, so the achievable memory bandwidth, and the number of passes over memory a circuit needs, decide the run time.

Data-centre FPGAs with high-bandwidth memory make this question concrete for reconfigurable hardware. The Alveo U55C exposes 16~GB of HBM2 as 32 independent pseudo-channels with a nominal aggregate bandwidth of about 460~GB/s~\cite{xilinx2023alveo}, in the same range as the GDDR6 of workstation GPUs. Whether an FPGA statevector engine can use that bandwidth, and where it then stands against CPU and GPU simulators, depends on design decisions that are invisible on a GPU: how the state is laid out over the channels, how many independent engines stream it, how wide and how deep their memory requests are, and how many gates are applied per pass.

This paper answers these questions with measurements. We start from an existing HBM statevector kernel for the U55C (the \emph{baseline}), show by measurement that it moves at most \NumVsevenBwMax~GB/s through HBM because it issues single-beat accesses, and build a redesigned simulator around streaming engines, on-chip tiles and a single multiplexed $2\times 2$ datapath primitive. We then measure the memory-system ceilings of the card, run a design-space study over engines, pseudo-channels, port width, tiling, fusion, layout and exchange, and compare the result with multi-threaded CPU simulators and a GPU on identical circuits, with an explicit timing boundary for every platform.

\subsection{Contributions}\label{subsec:contributions}

\begin{enumerate}
\item \textbf{FPGA architecture.} A statevector engine for HBM FPGAs in which every gate is executed by one multiplexed $2\times2$ primitive (MUX2), applied to $2^{12}$-amplitude tiles held in a parity-banked on-chip buffer at four amplitude pairs per cycle, with up to four high qubits gathered per tile. The engine runs as one compute unit per pseudo-channel, with dedicated exchange units for bank-qubit remaps. A single card simulates 30 qubits in complex64 (GHZ-30 in \NumFlagGhzThirtyEtoE~s end to end).
\item \textbf{Memory-system design-space study.} Measured HBM ceilings on the U55C (\NumHbmThirtytwoCopy~GB/s for 32 channels in copy mode, \NumHbmSixteenPerCURead~GB/s per channel), the scaling of the simulator with the number of engines and channels, and ablations of tile size, gate fusion, multi-gate passes, high-qubit slots, layout, port width and on-card versus host-mediated exchange.
\item \textbf{Software integration.} A host compiler that maps any Qiskit circuit to MUX2 records without decomposing one- and two-qubit gates, schedules passes and remaps, and a runtime that keeps state initialisation, exchange and sampling off the PCIe critical path.
\item \textbf{Functional correctness.} A fidelity-based test suite with generic entangled inputs on which the revised simulator passes all \NumSuiteTotal{} cases at $1-F\leq 10^{-5}$, including analytic checks at 30 qubits; an audit of the baseline that identifies an out-of-bounds CX loop on hardware.
\item \textbf{Cross-platform evaluation.} Measurements against Qiskit Aer, qsim and Qulacs on a 128-core CPU host and against an NVIDIA A100 GPU (Qiskit Aer with cuStateVec, and CUDA-Q) on the same circuits, with device-only and end-to-end timing, measured board power and energy, and latency distributions. One card is faster than the best CPU configuration for bandwidth-bound circuits from \NumXplatGhzCrossN{} qubits and slower, by a factor of \NumXplatCpuFasterOtherThirtyMin to \NumXplatCpuFasterOtherThirtyMax{} at 30 qubits, for compute-bound circuits. The A100 is faster than the card for compute-bound circuits and about as fast for GHZ at 30 qubits.
\end{enumerate}

The paper is organised as follows. Section~\ref{sec:related} reviews related work, Section~\ref{sec:baseline} describes the platform and the audit of the baseline kernel, Section~\ref{sec:architecture} presents the revised architecture and host software, Section~\ref{sec:methodology} the methodology, Section~\ref{sec:results} the results, Section~\ref{sec:discussion} discusses them, Section~\ref{sec:limitations} lists the limitations and Section~\ref{sec:conclusions} concludes.
